\def\baselinestretch{1}

\chapter{Methodology}\label{ch:methodology}

\def\baselinestretch{1.66}

\section{Introduction}\label{sec:methodology_introduction}

Trajectory generation is a well-studied problem. Trajectories with smooth higher-order derivatives are known to reduce the excitation of poorly damped structural modes. To evaluate the yield impact of our proposed trajectories without requiring physical scanner hardware, we need a simulation testbed where they can be executed and analyzed.

\section{Simulation}\label{sec:methodology_simulation}

In this section we describe the implementation of the dynamics model~\cite{Al-Rawashdeh2022_Caracterization} described in Chapter~\ref{ch:model} using the MuJoCo~\cite{todorov2012mujoco} physics engine for Python. The simulation serves as a testbed for validating the proposed trajectories and evaluating the synthetic wafer defect maps generated by their traversal.

The MuJoCo~\cite{todorov2012mujoco} Python bindings allow us to define the context of our simulation from an XML string by calling the \texttt{mujoco.MjModel.from\_xml\_string(MODEL\_XML)} method. We extend from MuJoCo's \texttt{<worldbody>} tag on which we place a root \texttt{<body>} tag, representing our \textit{base frame}; this allows us to organize the machine in the simulation using the three chain hierarchical tree. Table~\ref{tab:body_mapping} maps the specific names used in the MuJoCo XML to their counterparts in the structural diagram of Figure~\ref{fig:litho_model_full}.

\begin{table}[htbp]
\centering
\caption{Mapping of MuJoCo XML bodies to the model components}
\label{tab:body_mapping}
\begin{tabular}{ll}
\toprule
\texttt{<body>} \texttt{name} property & Model body \\
\midrule
\texttt{base\_frame}        & Base Frame ($\mathcal{F}_1$) \\
\texttt{w\_ls\_act}         & Wafer LS ($i_3=2$) \\
\texttt{w\_ls\_flex}        & flex coupling ($i_3=3$) \\
\texttt{w\_ss\_act}         & Wafer SS ($i_3=4$) \\
\texttt{w\_ss\_flex}        & flex coupling ($i_3=5$) \\
\texttt{w\_stage}           & Wafer Stage, passive ($i_3=6$) \\
\texttt{wafer}              & Die ($N_3=7$) \\
\texttt{r\_ls\_act}         & Reticle LS ($i_1=2$) \\
\texttt{r\_ls\_flex}        & flex coupling ($i_1=3$) \\
\texttt{r\_ss\_act}         & Reticle SS ($i_1=4$) \\
\texttt{r\_ss\_flex}        & flex coupling ($i_1=5$) \\
\texttt{r\_stage}           & Reticle Stage, passive ($i_1=6$) \\
\texttt{mask}               & Mask ($N_1=7$) \\
\texttt{metro\_frame}       & Metrology Frame ($i_2=2$) \\
\texttt{optics\_box}        & Optics Box ($i_2=3$) \\
\texttt{lens}               & Lens Element ($i_2=4$) \\
\bottomrule
\end{tabular}
\end{table}

The simulation environment is defined through several specialized XML attributes that handle different aspects of the physical world. Table~\ref{tab:mujoco_tags} summarizes the function of each tag within the context of our lithography scanner model.

\begin{table}[htbp]
\centering
\caption{MuJoCo XML tag functions in the scanner simulation}
\label{tab:mujoco_tags}
\begin{tabular}{ll}
\toprule
Tag & Function \\
\midrule
\texttt{<worldbody>} & Defines the inertial root, the floor plane ($\mathcal{F}_0$). \\
\texttt{<body>}      & Creates a rigid body with its own local coordinate system. \\
\texttt{<inertial>}  & Assigns mass $m_{i_j}$ and the aggregated inertia tensor $I_{i_j}$. \\
\texttt{<joint>}     & Implements the degrees of freedom $f_{i_j}$ and $\theta_{z_{i_j}}$. \\
\texttt{<geom>}      & Defines the material (granite, silicon, glass) and collision shapes. \\
\texttt{<actuator>}  & Models the Lorentz and moving-coil motors via servo controllers. \\
\bottomrule
\end{tabular}
\end{table}

The mechanical properties of the couplings, represented by the Kelvin-Voigt elements, are implemented via the \texttt{stiffness} and \texttt{damping} attributes of the \texttt{<joint>} tags. To begin, the common base frame ($\mathcal{F}_1$) is anchored to the world frame ($\mathcal{F}_0$) via a 3-DOF joint (two prismatic, one revolute) representing the isolation system with stiffness $K_1$ and damping $C_1$, as seen in Eq.~\eqref{eq:baseframeeq} and~\eqref{eq:baseframerot}. For the two actuated stages of each chain (long-stroke and short-stroke), where the physical stiffness is replaced by the controller's active response, the simulation uses high $K_P$ values ($\texttt{kp} = 10^7$\,N/m, to approximate the stiffness of a rigid servo motor while maintaining numerical stability with the implicit integrator at a timestep of $dt = 10^{-4}$\,s).

Tables~\ref{tab:param_matrix_wafer},~\ref{tab:param_matrix_reticle}, and~\ref{tab:param_matrix_optics} illustrate the distribution of these variables for the three chains.

\begin{table}[htbp]
\centering
\caption{Wafer chain stiffness ($K$) and damping ($C$) parameters}
\label{tab:param_matrix_wafer}
\begin{tabular}{lcc}
\toprule
Coupling Interface & Stiffness Variable ($K$) & Damping Variable ($C$) \\
\midrule
Floor $\rightarrow$ Base Frame   & \texttt{K1}       & \texttt{C1} \\
Base $\rightarrow$ LS Stage      & \texttt{KP\_ACT}  & \texttt{KV\_ACT\_TRANS} \\
LS Stage $\rightarrow$ LS Flex   & \texttt{K\_LS\_F} & \texttt{C\_LS\_F} \\
LS Flex $\rightarrow$ SS Stage   & \texttt{KP\_ACT}  & \texttt{KV\_ACT\_TRANS} \\
SS Stage $\rightarrow$ SS Flex   & \texttt{K\_SS\_F} & \texttt{C\_SS\_F} \\
SS Flex $\rightarrow$ Stage      & \texttt{K\_ST\_F} & \texttt{C\_ST\_F} \\
Stage $\rightarrow$ Die          & \texttt{K8\_W}    & \texttt{C8\_W} \\
\bottomrule
\end{tabular}
\end{table}

\begin{table}[htbp]
\centering
\caption{Reticle chain stiffness ($K$) and damping ($C$) parameters}
\label{tab:param_matrix_reticle}
\begin{tabular}{lcc}
\toprule
Coupling Interface & Stiffness Variable ($K$) & Damping Variable ($C$) \\
\midrule
Floor $\rightarrow$ Base Frame   & \texttt{K1}       & \texttt{C1} \\
Base $\rightarrow$ LS Stage      & \texttt{KP\_ACT}  & \texttt{KV\_ACT\_TRANS} \\
LS Stage $\rightarrow$ LS Flex   & \texttt{K\_LS\_F} & \texttt{C\_LS\_F} \\
LS Flex $\rightarrow$ SS Stage   & \texttt{KP\_ACT}  & \texttt{KV\_ACT\_TRANS} \\
SS Stage $\rightarrow$ SS Flex   & \texttt{K\_SS\_F} & \texttt{C\_SS\_F} \\
SS Flex $\rightarrow$ Stage      & \texttt{K\_ST\_F} & \texttt{C\_ST\_F} \\
Stage $\rightarrow$ Mask         & \texttt{K8\_R}    & \texttt{C8\_R} \\
\bottomrule
\end{tabular}
\end{table}

\begin{table}[htbp]
\centering
\caption{Optics chain stiffness ($K$) and damping ($C$) parameters}
\label{tab:param_matrix_optics}
\begin{tabular}{lcc}
\toprule
Coupling Interface & Stiffness Variable ($K$) & Damping Variable ($C$) \\
\midrule
Base Frame $\rightarrow$ Metro Frame & \texttt{K2\_O} & \texttt{C2\_O} \\
Metro Frame $\rightarrow$ Optics Box & \texttt{K3\_O} & \texttt{C3\_O} \\
Optics Box $\rightarrow$ Lens Element & \texttt{K4\_O} & \texttt{C4\_O} \\
\bottomrule
\end{tabular}
\end{table}

\begin{figure}[htbp]
    \centering
    \includegraphics[width=\textwidth]{scanning_3d.png}
    \caption{Screenshot of the MuJoCo simulation in the midst of the scanning motion, showing the wafer and the optics chain under the metrology frame}
    \label{fig:scanning_simulation}
\end{figure}

To successfully adapt our planar reduction of the reference model~\cite{Al-Rawashdeh2022_Caracterization} into MuJoCo's multibody solver, two numerical adaptations are in order: 

\begin{enumerate}
    \item Native joint damping against absolute velocity would exert artificial drag ($K_V/K_P = 10$\,mm per m/s) during cruise, so we apply the velocity gain $K_V$ to the tracking-error velocity $(\dot{r} - \dot{y})$ within the control loop to represent the frictionless air bearings.

    \item Actuators transmit reaction forces to the parent body by default, so we apply reaction cancellation to the base frame at each integration step to represent balance masses used to absorb long-stroke reaction forces~\cite{Butler2011_PositionControl}, which the analytical model~\cite{Al-Rawashdeh2022_Caracterization} represents as external stage forces.
\end{enumerate}

Synchronized scanning motion is commanded with nominal cruise velocities of $v_{w,\mathrm{scan}} = 0.8$\,m/s and $v_{r,\mathrm{scan}} = 3.2$\,m/s. Under acceleration transients up to $a_{w,\max} = 20$\,m/s$^2$, this setup captures the full-state dynamic interactions derived in Section~\ref{sec:model_dynamics}.

Baseline control configurations (such as Case~1 in~\cite{Al-Rawashdeh2022_Caracterization}) are realized on this plant by disabling active fine-stage feedback loops.

\section{Reference Trajectory Generation}\label{sec:methodology_trajectories}

The trajectories commanded to the stages are \textbf{fourth-order profiles}: piecewise-polynomial motions in which position, velocity, acceleration, and jerk are all continuous, and the fourth derivative --- the \textit{snap} --- switches between $\pm s_{\max}$ and zero~\cite{Lambrechts2003_Trajectories}. Within each segment the position evolves as
\begin{equation}
    p(t) = p_0 + v_0 t + \frac{a_0}{2}t^2 + \frac{j_0}{6}t^3 + \frac{s}{24}t^4,
    \label{eq:fourth_order_poly}
\end{equation}
where ($p_0$, $v_0$, $a_0$, $j_0$) are the boundary conditions inherited from the previous segment and $s \in \{-s_{\max}, 0, +s_{\max}\}$. A symmetric point-to-point motion consists of fifteen such segments, arranged so that none of the kinematic bounds ($v_{\max}$, $a_{\max}$, $j_{\max}$, $s_{\max}$) is exceeded.

We implement an analytical trajectory generator that structures the motion profile across fifteen segments (seven per acceleration ramp) where snap is held constant. Physical limitations are addressed by the following considerations from~\cite{Lambrechts2003_Trajectories}:
\begin{itemize}
    \item When the jerk plateau cannot be reached, the effective jerk reduces to $\sqrt{a_{\max} s_{\max}}$
    \item If the acceleration limit is unreachable, then the effective acceleration follows from the corresponding quadratic equation; if the jerk phase collapses as well, the snap-phase duration instead follows from a cube-root solve
    \item If the total travel distance is too short to accommodate the required transition ramps, then the cruise velocity is iteratively reduced via bisection.
\end{itemize}
To streamline the trials, this same algorithmic structure is used to derive the lower order profiles: i.e., a second-order, acceleration-bounded trapezoidal profile and a third-order, seven-segment jerk-bounded S-curve.

The order of the motion profile determines the high-frequency content of the force commanded to initiate such motion: a second-order profile demands instantaneous jumps in acceleration, while a third-order profile is discontinuous in jerk~\cite{Dong2023_SCurve}, so both inject energy across a broad frequency band. A fourth-order trajectory, whose snap is bounded as noted above, is reported to roll off the excitation spectrum at $-80$\,dB per decade~\cite{Lambrechts2003_Trajectories} --- a theoretical asymptotic for constant-snap profiles, not separately verified spectrally in this work --- reducing the response of poorly damped structural modes such as those of the end-effectors, the reticle and the wafer. The settling period after each step~\cite{Lambrechts2003_Trajectories, Al-Rawashdeh2023_Trajectories} is thereby minimized, which is essential for the constant-velocity scanning phase, where ideally no force beyond the viscous terms is present.

\section{Stage Feedback Control}\label{sec:methodology_control}

The stage controllers are evaluated across four incremental configurations, following~\cite{Al-Rawashdeh2022_Caracterization}. Each setup introduces a specific feedforward or feedback component to isolate its impact on exposure performance in Chapter~\ref{ch:experiments}.

% (collocated long-stroke control)
In Case 1, each long-stroke stage is regulated in closed loop by a collocated PD controller acting on the long-stroke encoder $y_c(t)$:
\begin{equation}
    u(t) = K_P\, e(t) + K_V\, \dot{e}(t),
    \label{eq:case1_servo}
\end{equation}
where $e(t) = r(t) - y_c(t)$ is the long-stroke tracking error, with gains $K_P = 10^7$\,N/m and $K_V = 10^5$\,N$\cdot$s/m. Position and velocity are read from the long-stroke stage rather than the end-effector because non-collocated control through the lightly damped reticle coupling ($\zeta \approx 0.04$) destabilizes the loop. Consequently, the lag between the long-stroke stage and the end-effector remains uncompensated and appears directly in the synchronization error. On this frictionless plant, this lag is dominated by its acceleration-proportional component: servo deflection $m\,a/K_P$ and quasi-static coupling stretch $m_d\,a/K$ carrying the downstream mass $m_d$. This configuration serves as the baseline validated in Experiment~1 (Section~\ref{sec:exp_case1}).

% (collocated long-stroke control)
Since both lag components are functions of the known reference kinematics, in Case 2 they are cancelled by commanding the long stroke \textit{ahead} of the reference:
\begin{equation}
    r_{LS}(t) = r(t) - \big(\ell_v\, \dot{r}(t) + \ell_a\, \ddot{r}(t)\big),
    \label{eq:lag_ff}
\end{equation}
so that the end-effector may land on the desired position. The coefficients ($\ell_v$, $\ell_a$) for each chain are identified by performing iterative least-squares regressions on the end-effector lag with respect to reference velocity and acceleration. Each resulting fit is added to the feedforward control loop while subsequent passes repeat the regression on the remaining residuals. Three iterations are required to prevent post-acceleration coupling vibrations from skewing the quasi-static fit.This is a model-free equivalent to the spring-compensation configuration of the reference study~\cite{Al-Rawashdeh2022_Caracterization}.
%, and requires no additional hardware or sensing. What the fit cannot represent --- and what remains after Case 2 --- is precisely that dynamic transient, which motivates the fine-stage feedback below. This is the model-free counterpart of the spring-compensation configuration of the reference study, and requires no additional hardware or sensing.

% (fine-stage feedback)
The error remaining after feedforward is a non-repeatable transient response that can only be suppressed through feedback. Thus, in Case 3, the short-stroke stage of each chain implements a closed position loop on the end-effector error $e_f(t)$, commanding the short-stroke actuator directly. In actual step-and-scan systems, this measurement is readily available from interferometry, as required for exposure metrology. Case 3a uses a linear PI law with anti-windup, whereas Case 3b replaces the linear integrator with a Hybrid Integrator--Gain System (HIGS)~\cite{Heertjes2025_NonlinearControl}, given by:
\begin{equation}
    \dot{x}_1 = \omega_h\, e_f \quad \text{while } e_f\, x_1 \geq x_1^2 / k_h,
    \qquad x_1 = k_h\, e_f \quad \text{otherwise},
    \label{eq:higs}
\end{equation}
with output $u_f = k_p e_f + \omega_i \int x_1\, dt$. By resetting its state to the gain line $k_h e_f$ whenever the sector condition is violated, the HIGS element provides integral action with roughly $38^\circ$ less phase lag than a linear integrator, relaxing the overshoot-versus-stiffness limitation of linear designs --- the property that motivates nonlinear integrators for wafer stages. Both fine loops saturate at the physical travel of the short stroke, with conditional integration preventing windup during large transients. The loop gains are tuned empirically over a grid up to the stability boundary of the simulated plant: the PI loop uses $k_p = 16$, $k_i = 480$\,s$^{-1}$, and the HIGS loop $k_p = 16$, $\omega_i = 2\pi \cdot 200$\,rad/s, $\omega_h = 2\pi \cdot 80$\,rad/s, $k_h = 1$; one step further in gain destabilizes both loops against the $180$\,Hz flexible-chain modes. The HIGS element sustains an integrator frequency roughly $42\times$ above the PI's equivalent $k_i/k_p$ before reaching that boundary --- the phase advantage of the nonlinear integrator expressed as usable gain. The fine-stage authority is ultimately bounded by the short-stroke actuator's own velocity loop (pole at $K_P/K_V = 100$\,rad/s), a structural property of this generic replication that Chapter~\ref{ch:experiments} quantifies.

\section{Population-Based Trajectory Search}\label{sec:methodology_pso}

The optimization problem~\eqref{eq:problem} is searched with particle swarm optimization (PSO) over the four kinematic bounds. Each particle holds a position $\mathbf{x} = (v_{\max}, a_{\max}, j_{\max}, s_{\max})$ and is updated by the standard rule
\begin{equation}
    \mathbf{u}_{k+1} = \omega\, \mathbf{u}_k + c_1 r_1 (\mathbf{p}_k - \mathbf{x}_k) + c_2 r_2 (\mathbf{g}_k - \mathbf{x}_k), \qquad \mathbf{x}_{k+1} = \mathbf{x}_k + \mathbf{u}_{k+1},
    \label{eq:pso}
\end{equation}
with inertia $\omega = 0.6$, cognitive and social weights $c_1 = c_2 = 1.4$, velocities clamped to $\pm 0.5$ of the search-box span per dimension, and positions clipped to the search box. Evaluating one particle runs the full step-and-scan sequence over the wafer under the candidate trajectory, converts the per-die exposure indices into a wafer map, and prices the map. Two cost definitions are compared as the scenarios E3 and E4 of the evaluation plan: a conventional servo cost (the wafer-average exposure MSD, normalized by the failure threshold) and the yield-aware cost built on the CNN. For the latter, the weighted form~\eqref{eq:cost} is instantiated as
\begin{equation}
    \mathcal{J}_{map} = \big(1 - P_{none}\big) + \sum_{c\, \in\, \mathcal{C}_{ood}} P_c,
    \label{eq:cost_corrected}
\end{equation}
where $\mathcal{C}_{ood}$ is the set of classes that scanner motion cannot plausibly produce. The first term prices any departure from a clean map, so that the optimizer can spend an \textit{acceptable-pattern} yield loss when it is worth it; the second is a hard wall that keeps it from cheapening the wafer into a morphology the classifier attributes to non-scanner mechanisms. Both scenarios share a throughput reward $\mathcal{J} = \mathcal{J}_{map} + \gamma\, T_{wafer}/T_{ref}$ with $\gamma = 1$ and $T_{ref} = 15$\,s, so that a faster wafer is cheaper and the search trades cycle time against the map cost rather than merely respecting a budget; $T_{wafer}$ follows from the scan-profile duration, the settling budget, and the die count.

\section{From Synchronization Error to Wafer Maps}\label{sec:methodology_cnn}

The link between trajectory parameters and yield is established in three steps: the simulated synchronization error is aggregated per die, the resulting pass/fail pattern forms a simulated wafer map, and a convolutional neural network classifies that map into a defect-morphology class.

The CNN classifier is trained on WM-811K~\cite{Wu2015_WM811K}, a public dataset of 811{,}457 wafer maps collected from real fabrication lines, of which 172{,}950 carry a manually assigned failure label. Each map is a grid of dies marked as pass or fail, annotated with one of nine morphology classes: \textit{none}, \textit{Center}, \textit{Donut}, \textit{Edge-Loc}, \textit{Edge-Ring}, \textit{Loc}, \textit{Near-full}, \textit{Random}, and \textit{Scratch}. The labelled maps are split 80/20 with class stratification, and every map is resized to $64 \times 64$ cells to give the network a uniform input.

The network architecture comprises three convolutional blocks (16, 32, and 64 filters of size $3\times3$, each followed by ReLU activation and $2\times2$ max-pooling) and two fully connected layers (128 hidden units with dropout of 0.5, and a 9-way softmax output). It is trained for 5 epochs with the Adam optimizer (learning rate $10^{-3}$, batch size 32) under a cross-entropy loss, retaining the weights with the highest hold-out accuracy.

To generate simulated maps, for each die scanned in simulation, the synchronization error recorded during its exposure is reduced to a scalar criterion --- either the maximum tracking error or the MA/MSD pair of Chapter~\ref{ch:control_objectives} --- and compared against a threshold to declare the die good or failing. The resulting grid, arranged in the same encoding as WM-811K (background, pass, fail), is resized to $64\times64$ and presented to the CNN, which returns the probability distribution $P_c$ used by the cost function~\eqref{eq:cost}.

For validation, a physically meaningful trajectory-to-defect mapping must activate only the classes plausibly caused by scanner motion: \textit{none}, \textit{Scratch}, \textit{Edge-Loc}, and \textit{Random}. The remaining classes correspond to etching, cleaning, and deposition signatures; if simulated maps systematically activated them, the mapping --- not the machine --- would be at fault. This criterion, together with the requirement that the cost shift reproducibly under sweeps of $(v_{\max}, a_{\max}, j_{\max}, s_{\max})$, defines what constitutes a successful validation in Chapter~\ref{ch:experiments}.

